\documentclass[11pt]{article}
\usepackage[a4paper,margin=30mm]{geometry}
\usepackage{amsmath,amssymb,amsthm,mathtools}
\usepackage[hidelinks]{hyperref}

\newtheorem{proposition}{Proposition}
\newtheorem{remark}{Remark}
\newcommand{\Hh}[2]{H_{#1}^{(#2)}}

\title{The Laurent--Jet Origin of the Compact Brown--Zudilin Weights}
\author{}
\date{}

\begin{document}
\maketitle

\section{The compact weights}

Put
\[
 T_{n,k,\ell}
 =
 \binom{n+k}{n}\binom nk^2
 \binom{n+\ell}{n}\binom n\ell^2
 \binom{n+k+\ell}{n}.
\]
For \(r\geq1\), write
\[
 A_r(x)=\Hh{n+x}{r}-\Hh{x}{r},
 \qquad
 B_r(x)=\Hh{n-x}{r}-\Hh{x}{r},
\]
and define
\[
 \alpha=A_1(k)-A_1(\ell),\qquad
 \beta=B_1(k)-B_1(\ell),\qquad
 \Psi=\frac{\alpha}{2}+\beta.
\]
The experimentally discovered weights are
\[
 \widehat w_3
 =\Hh{n+k}{3}-\Psi\Hh{n+k}{2}
\]
and
\[
\begin{split}
 w_5={}&\Hh{n+k}{5}
 +\frac{\alpha-\beta}{2}\Hh{n+k}{4}\\
 &+
 \left\{
 \frac{A_2(k)+A_2(\ell)}4
 -\frac{\alpha\Psi}{2}
 \right\}\Hh{n+k}{3}.
\end{split}
\]
We explain why precisely these combinations occur.

\section{The rational Barnes kernel}

In the totally symmetric Brown--Zudilin Barnes integral, application of
\(\Gamma(z)\Gamma(1-z)=\pi/\sin(\pi z)\) separates the integrand into three
sine factors and the rational function
\[
 \mathcal R_n(s,t)=R_n(s)R_n(t)C_n(s+t),
\]
where
\[
 R_n(s)=
 \frac{\prod_{i=1}^{n}(s+i)}
      {\prod_{j=n+1}^{2n+1}(s+j)^2},
 \qquad
 C_n(u)=\prod_{j=n+2}^{2n+1}(u+j).
\]
The finite double poles are
\[
 s_k=-n-k-1,\qquad t_\ell=-n-\ell-1,
 \qquad 0\leq k,\ell\leq n.
\]

\begin{proposition}[Leading Laurent coefficient]
\label{prop:leading}
At \((s_k,t_\ell)\),
\[
 \lim_{s\to s_k,\ t\to t_\ell}
 (s-s_k)^2(t-t_\ell)^2\mathcal R_n(s,t)
 =
 \frac{(-1)^n}{n!}\,T_{n,k,\ell}.
\]
\end{proposition}

\begin{proof}
First,
\[
 \lim_{s\to s_k}(s-s_k)^2R_n(s)
 =
 (-1)^n\frac{(n+k)!}{k!^3(n-k)!^2}.
\]
Indeed, the numerator at \(s_k\) is
\((-1)^n(n+k)!/k!\), whereas the product of the remaining denominator
factors, before squaring, is
\((-1)^k k!(n-k)!\).
Similarly for \(t_\ell\).  Finally,
\[
 C_n(s_k+t_\ell)
 =
 (-1)^n\frac{(n+k+\ell)!}{(k+\ell)!}.
\]
Multiplication gives
\[
 (-1)^n
 \frac{(n+k)!(n+\ell)!(n+k+\ell)!}
 {k!^3\ell!^3(k+\ell)!(n-k)!^2(n-\ell)!^2}
 =\frac{(-1)^n}{n!}T_{n,k,\ell}.
\qedhere
\]
\end{proof}

Thus the binomial kernel itself is forced: it is the leading coefficient of
the finite double-pole Laurent expansion.

\section{The anti-diagonal jet}

Let
\[
 \Lambda_{k,\ell}
 =
 \frac{(-1)^n}{n!}T_{n,k,\ell}
\]
and normalize the local germ by
\[
 E_{k,\ell}(x,y)=
 \frac{x^2y^2
 \mathcal R_n(s_k+x,t_\ell+y)}
 {\Lambda_{k,\ell}}.
\]
Then \(E_{k,\ell}(0,0)=1\).
The decisive direction is the anti-diagonal derivative
\[
 \nabla=\partial_x-\partial_y.
\]
Since \(C_n\) depends only upon \(s+t\), it is constant on
\((s,t)=(s_k+x,t_\ell-x)\).  Consequently \(\nabla\) annihilates the entire
coupling factor \(C_n(s+t)\).

To separate the two elementary pieces of \(R_n\), put
\[
 N_n(s)=\prod_{i=1}^n(s+i)
\]
and, locally at \(s_k\),
\[
 D_{n,k}(s)=
 \prod_{\substack{j=n+1\\j\neq n+k+1}}^{2n+1}(s+j).
\]
Define the normalized anti-diagonal germs
\[
 \mathcal N_{k,\ell}(x)=
 \frac{N_n(s_k+x)N_n(t_\ell-x)}
      {N_n(s_k)N_n(t_\ell)}
\]
and
\[
 \mathcal D_{k,\ell}(x)=
 \frac{D_{n,k}(s_k+x)D_{n,\ell}(t_\ell-x)}
      {D_{n,k}(s_k)D_{n,\ell}(t_\ell)}.
\]
Finally set
\[
 e_1=\left.\frac{d}{dx}
 \log E_{k,\ell}(x,-x)\right|_{x=0},
\quad
 \nu_j=\left.\frac{d^j}{dx^j}
 \log\mathcal N_{k,\ell}(x)\right|_{x=0},
\quad
 \delta_1=\left.\frac{d}{dx}
 \log\mathcal D_{k,\ell}(x)\right|_{x=0}.
\]

\begin{proposition}[Harmonic coordinates are Laurent jets]
\label{prop:jets}
One has
\[
 \nu_1=-\alpha,\qquad
 \delta_1=\beta,\qquad
 e_1=-2\Psi,\qquad
 \nu_2=-\bigl(A_2(k)+A_2(\ell)\bigr).
\]
\end{proposition}

\begin{proof}
At \(s=s_k\),
\[
 \frac{d}{ds}\log N_n(s)
 =
 \sum_{i=1}^n\frac1{s+i}
 =
 -A_1(k).
\]
The sign of the \(\ell\)-term is reversed by \(t=t_\ell-x\); hence
\(\nu_1=-A_1(k)+A_1(\ell)=-\alpha\).

The nonvanishing denominator differences at \(s_k\) are
\[
 -k,\ldots,-1,1,\ldots,n-k,
\]
and therefore
\[
 \left.\frac{d}{ds}\log D_{n,k}(s)\right|_{s=s_k}
 =-H_k+H_{n-k}=B_1(k).
\]
The anti-diagonal difference gives
\(\delta_1=B_1(k)-B_1(\ell)=\beta\).

Locally, apart from the double-pole monomials,
\[
 \log E=
 \log N_n(s)+\log N_n(t)
 -2\log D_{n,k}(s)-2\log D_{n,\ell}(t)
 +\log C_n(s+t)+\text{constant}.
\]
The last term is killed in the anti-diagonal direction.  Thus
\[
 e_1=-\alpha-2\beta=-2\Psi.
\]

Finally,
\[
 \left.\frac{d^2}{ds^2}\log N_n(s)\right|_{s=s_k}
 =
 -\sum_{i=1}^n\frac1{(s_k+i)^2}
 =-A_2(k).
\]
The second derivative is unchanged by replacing \(x\) with \(-x\), so
\(\nu_2=-A_2(k)-A_2(\ell)\).
\end{proof}

\section{Reconstruction of the two compact weights}

Put \(h_r=\Hh{n+k}{r}\).  Proposition~\ref{prop:jets} gives the identities
\[
 \boxed{\quad
 \widehat w_3=h_3+\frac{e_1}{2}h_2
 \quad}
\]
and
\[
 \boxed{\quad
 w_5
 =h_5-\frac{\nu_1+\delta_1}{2}h_4
 -\frac{\nu_2+\nu_1e_1}{4}h_3.
 \quad}
\]
Indeed,
\[
 \frac{e_1}{2}=-\Psi,\qquad
 -\frac{\nu_1+\delta_1}{2}
 =\frac{\alpha-\beta}{2},
\]
and
\[
 -\frac{\nu_2+\nu_1e_1}{4}
 =
 \frac{A_2(k)+A_2(\ell)}4-\frac{\alpha\Psi}{2}.
\]

This proves that every ingredient of the compact weights has a canonical
local meaning:
\[
\begin{array}{c|c}
\text{quantity}&\text{Laurent--jet origin}\\ \hline
T_{n,k,\ell}&\text{leading double-pole coefficient}\\
\alpha&-\nabla\log(\text{numerator})\\
\beta&\nabla\log(\text{regular denominator})\\
\Psi&-\frac12\nabla\log(\text{whole local germ})\\
A_2(k)+A_2(\ell)&-\nabla^2\log(\text{numerator}).
\end{array}
\]
Most importantly, the direction \(\partial_s-\partial_t\) removes the coupled
factor \(C_n(s+t)\).  This explains the otherwise surprising absence of
harmonic numbers at the argument \(n+k+\ell\) from the final formulas.

\section{What remains for a complete proof of the companion identities}

The preceding propositions rigorously explain the kernel and the compact
weights, but one further analytic lemma is required to identify their sums
with the Brown--Zudilin coefficients \(P_n\) and \(\widehat P_n\).
One must carry the partial-fraction decomposition of \(\mathcal R_n(s,t)\)
through the universal sine-kernel integrals
\[
 I_{a,b}^{(r,s)}
 =
 \frac1{(2\pi i)^2}
 \iint
 \frac{1}{(u+a)^r(v+b)^s}
 \frac{\pi}{\sin\pi u}
 \frac{\pi}{\sin\pi v}
 \frac{\pi}{\sin\pi(u+v)}
 \,du\,dv,
\qquad r,s\in\{1,2\}.
\]
The required universal coefficient-extraction statement is that the
\(\zeta(2)\)-rational and rational weight-five projections of these four
integrals pair with the local Laurent coefficients exactly as
\[
 h_3+\frac{e_1}{2}h_2
\]
and
\[
 h_5-\frac{\nu_1+\delta_1}{2}h_4
 -\frac{\nu_2+\nu_1e_1}{4}h_3.
\]
Once this universal lemma is written and proved, summing
Proposition~\ref{prop:leading} over the finite double poles gives immediately
\[
 \widehat P_n=\sum_{k,\ell=0}^nT_{n,k,\ell}\widehat w_3(n,k,\ell),
\qquad
 P_n=\sum_{k,\ell=0}^nT_{n,k,\ell}w_5(n,k,\ell).
\]

\begin{remark}
Brown and Zudilin explicitly describe this route: rewrite the Barnes
integrand as a rational function times reciprocal sine factors, decompose the
rational part into partial fractions, and evaluate the resulting universal
double integrals.  The calculation above identifies the previously missing
local structure of that decomposition.  It also gives a systematic discovery
method for higher cellular families: compute finite-pole Laurent coefficients,
then take jets in directions tangent to the level sets of the coupled Gamma
arguments.
\end{remark}

\end{document}
